\section{Vertical Order Traversal}
\label{sec:trees-vertical}

\problemstatement{Assign the root a horizontal distance (HD) of $0$; a left
child has HD one less than its parent, a right child one more. Return
nodes grouped by HD from leftmost to rightmost. Within the same HD, order
by depth (top to bottom); nodes at the same HD \emph{and} same depth are
ordered by value.}

\begin{patternbox}
\emph{``Vertical,'' ``columns,'' ``horizontal distance''} says: tag each
node with a \textbf{(HD, depth)} coordinate, collect everything into a map
keyed by HD, and sort within each column. The traversal is ordinary BFS or
DFS; the work is in the coordinate bookkeeping and the tie-breaking rules.
\end{patternbox}

\subsection*{Understanding the problem carefully}

Imagine dropping vertical lines through the tree. Every node falls on one
line, identified by its horizontal distance. Reading the answer means
listing lines left to right, and within a line, top to bottom. The precise
LeetCode-style rule adds a final tie-break: if two nodes share both HD and
depth, the smaller value comes first. Storing each node as a triple
$(\text{HD},\text{depth},\text{value})$ and sorting makes all three rules
fall out of one comparison.

\begin{ideabox}[Order by (HD, then depth, then value)]
Traverse the tree recording, for each node, the triple $(HD, depth,
val)$: root is $(0,0,val)$, a left child inherits $(HD-1, depth+1)$, a
right child $(HD+1, depth+1)$. Group into a map from HD to the list of
$(depth, val)$ pairs; sort each list by $(depth, val)$; emit columns in
increasing HD.
\end{ideabox}

\subsection*{Algorithm}

\begin{enumerate}
  \item BFS/DFS from the root carrying $(HD, depth)$; append $(depth,
        val)$ into $\textit{cols}[HD]$.
  \item For each HD in increasing order, sort its list by $(depth, val)$
        and output the values.
\end{enumerate}

\subsection*{Dry run}

\dryrun{Tree: root $1$; left $2$; right $3$; $2$'s right child $4$; $3$'s
left child $5$ (so $4$ and $5$ both land on HD $0$).}

\begin{itemize}
  \item $1:(0,0)$; $2:(-1,1)$; $3:(1,1)$; $4:(0,2)$; $5:(0,2)$.
  \item HD $-1$: $[2]$. HD $0$: $\{(0,1),(2,4),(2,5)\}$ sorted by
        (depth,val) $\to 1,4,5$. HD $1$: $[3]$.
  \item Result: $[[2],[1,4,5],[3]]$.
\end{itemize}

\subsection*{C++ implementation}

\begin{lstlisting}
#include <bits/stdc++.h>
using namespace std;

struct TreeNode {
    int val;
    TreeNode* left;
    TreeNode* right;
    TreeNode(int v) : val(v), left(nullptr), right(nullptr) {}
};

void dfs(TreeNode* node, int hd, int depth,
         map<int, vector<pair<int,int>>>& cols) {
    if (node == nullptr) return;
    cols[hd].push_back({depth, node->val});
    dfs(node->left,  hd - 1, depth + 1, cols);
    dfs(node->right, hd + 1, depth + 1, cols);
}

vector<vector<int>> verticalTraversal(TreeNode* root) {
    map<int, vector<pair<int,int>>> cols;   // hd -> (depth, val)
    dfs(root, 0, 0, cols);

    vector<vector<int>> ans;
    for (auto& [hd, vec] : cols) {          // map iterates hd ascending
        sort(vec.begin(), vec.end());       // by depth, then value
        vector<int> col;
        for (auto& [d, v] : vec) col.push_back(v);
        ans.push_back(col);
    }
    return ans;
}

int main() {
    TreeNode* root = new TreeNode(1);
    root->left = new TreeNode(2);
    root->right = new TreeNode(3);
    root->left->right = new TreeNode(4);
    root->right->left = new TreeNode(5);

    for (auto& col : verticalTraversal(root)) {
        for (int x : col) cout << x << " ";
        cout << "\n";
    }
    return 0;
}
\end{lstlisting}

\begin{complexitybox}
\textbf{Time Complexity.} $O(n\log n)$ -- dominated by sorting the columns
(total $n$ entries across all columns). \textbf{Space Complexity.} $O(n)$
for the map and the recursion stack.
\end{complexitybox}

\begin{takeawaybox}
Tagging nodes with coordinates and offloading the ordering to a sorted map
plus a per-column sort turns a geometric-sounding problem into
bookkeeping. Top view and bottom view (next sections) reuse the same HD
tag with a simpler rule for which node on each line ``wins.''
\end{takeawaybox}
